% !TeX root = ../main.tex
\begin{chineseabstract}
    % \footnotetext{*本研究得到某某基金（编号：）的资助}
    % \astfootnote{本研究得到某某基金（编号：）的资助}
    % \zhlipsum[2,5]
    
    随着智能穿戴设备与无线感知技术的普及，基于惯性测量单元（IMU）与WiFi的协同感知已成为人体健康监测领域的研究热点。IMU与WiFi具有成本低、易部署及隐私保护强等天然优势，但在实际应用中，单一传感模态在复杂环境下的特征表征能力受限。同时，受试者个体的生理特征与运动习惯存在显著差异，由此产生的领域偏移会导致预训练模型在未见受试者上出现严重的性能退化。此外，现有的人体姿态重构研究大多聚焦于当前时刻的静态还原，缺乏对未来趋势的动态预测能力。针对异构数据深度融合难、个体差异泛化瓶颈以及姿态重构缺乏预测性这三个核心问题，本文分别开展了基于多模态融合的人体活动识别方法及基于自回归的人体姿态重构与预测方法研究。

    在人体活动识别任务中，针对IMU与WiFi数据物理属性差异大、特征难以对齐和深度融合的问题，本文提出了一种基于多模态融合的人体活动识别方法（WICTIVITY）。该方法采用跨模态对比学习策略，显式拉近IMU的动力学特征和WiFi的空间场特征在潜在语义空间中的距离，并结合交叉注意力机制实现特征的深度互补。实验结果表明，WICTIVITY在XRFV2数据集上能够有效挖掘多模态互补信息，活动识别准确率达到96.06\%，身份与场景识别准确率分别达到99.93\%和99.99\%。针对跨受试者活动识别性能退化的问题，本文进一步提出一种基于全局语义锚点对齐的少样本适配策略。该策略通过引入轻量化瓶颈适配器，并利用原型对齐损失函数引导目标受试者的特征分布向源域习得的各类活动特征中心靠拢，从而有效校准个体差异引发的特征偏移。实验结果显示，模型在未见受试者上的平均初始识别率为72.57\%，仅通过5至10个目标域样本微调，平均准确率即可提升至78.82\%，并随样本量增加最终稳定在90.99\%。

    在人体姿态重构任务中，针对现有方法缺乏动态预测能力的问题，本文提出了一种基于自回归机制的人体姿态重构与预测方法（AIPose），旨在将人体姿态重构从静态映射扩展为动态时序生成问题。该方法构建了包含IMU特征提取、姿态重构及IMU预测的闭环生成式架构。通过将预测的未来时刻的IMU数据重新反馈至IMU特征提取器，实现了对人体姿态的连续重构与超前预测。实验结果表明，AIPose在XRFV2数据集上可以实现厘米级的姿态重构，并具备稳定的长时序姿态预测能力，其对当前时刻姿态重构的平均关节位置误差为49.35毫米，对未来时刻姿态的预测误差为55.19毫米。

    % 本文做出了以下贡献：
    
    % \begin{enumerate}[wide,]
    %     \item \zhlipsum[1]
    %     \item \zhlipsum[3]
    % \end{enumerate}
    
    \chinesekeywordstype{姿态重构；活动识别；惯性测量单元；WiFi；多模态融合；自回归}{应用研究}
    
\end{chineseabstract}